% ============================================================
\section{Programmazione Dinamica}
% Slide: s09_1_programmazioneDinamica-handout.pdf
% Fonte: Capitolo 14 CLRS — Prof. Roberto Posenato, A.A. 2025/2026
% ============================================================

\begin{softnote}
La \textbf{programmazione dinamica} è un paradigma per la progettazione di algoritmi
che risolvono problemi di ottimizzazione. Introdotta da Bellman (1957), si applica
quando il problema presenta \emph{sottostruttura ottima} e \emph{sottoproblemi
	sovrapposti}: la soluzione ottima contiene le soluzioni ottime dei sottoproblemi, e
un approccio ricorsivo naïve risolverebbe gli stessi sottoproblemi più volte,
con costo esponenziale.
\end{softnote}

% -------------------------------------------------------
\subsection{Idea generale}
% -------------------------------------------------------

La tecnica si articola in \textbf{4 fasi}:

\begin{enumerate}
\item \textbf{Caratterizzazione} della struttura di una soluzione ottima.
\item \textbf{Definizione ricorsiva} del valore di una soluzione ottima.
\item \textbf{Calcolo} del valore secondo uno schema \emph{bottom-up}.
\item \textbf{Costruzione} della soluzione ottima dalle informazioni calcolate.
\end{enumerate}

In termini moderni, la programmazione dinamica descrive il problema come un insieme
di \textbf{stati} (o sottoproblemi). Per ogni stato si definiscono:
\begin{itemize}
\item i \emph{casi base};
\item le \emph{transizioni}: come combinare stati più piccoli per ottenere il
valore dello stato corrente;
\item l'\emph{ordine di calcolo}, che deve rispettare le dipendenze tra stati.
\end{itemize}

Il costo totale dell'algoritmo è spesso:
\[
\#\text{stati} \;\times\; \text{costo per calcolare una transizione.}
\]

La stessa ricorrenza può essere implementata con una tabella \emph{bottom-up}
(tabulazione) oppure con ricorsione e \emph{memoization}.

\begin{nota}
\paragraph{Quando si usa la programmazione dinamica.}
Si usa quando:
\begin{itemize}
	\item una soluzione ottima contiene al suo interno le soluzioni ottime dei
	sottoproblemi (\textbf{sottostruttura ottima});
	\item un algoritmo ricorsivo richiederebbe di risolvere più volte gli stessi
	sottoproblemi (\textbf{sottoproblemi sovrapposti});
	\item i sottoproblemi possono essere descritti da un numero limitato di parametri,
	cioè da uno \textbf{stato}.
\end{itemize}
\end{nota}

\subsubsection{Confronto con Divide et Impera}

\begin{center}
\begin{tabular}{lll}
	\toprule
	\textbf{Proprietà} & \textbf{Tabulazione (bottom-up)} & \textbf{Divide et Impera} \\
	\midrule
	Tipo di risoluzione   & bottom-up   & top-down \\
	Tipo di programma     & iterativa   & ricorsiva \\
	Uso della memoria     & esplicito per i risultati & implicito \\
	Applicazione ideale   & sottoistanze non disgiunte & sottoistanze disgiunte \\
	\bottomrule
\end{tabular}
\end{center}

\begin{nota}
La programmazione dinamica può essere implementata anche top-down tramite
\emph{memoization}, mantenendo la struttura ricorsiva del divide et impera ma
riutilizzando i risultati già calcolati.
\end{nota}

% -------------------------------------------------------
\subsection{Esempi di applicazione: confronto di sequenze DNA}
% -------------------------------------------------------

\begin{softnote}
Una sequenza DNA si rappresenta come una stringa sull'alfabeto $\{A, C, G, T\}$.
Il confronto tra due sequenze può rivelare quanto siano correlati i due organismi
di origine. Esistono diversi criteri di confronto:
\begin{itemize}
	\item verificare se una sequenza è sottostringa dell'altra;
	\item determinare il numero minimo di modifiche per trasformare una nell'altra
	(distanza di edit);
	\item trovare una \emph{sottosequenza comune} di lunghezza massima.
\end{itemize}
\end{softnote}

% -------------------------------------------------------
\subsection{Massima Sottosequenza Comune (LCS)}
% -------------------------------------------------------

\begin{definizione}[Sottosequenza]
Data una sequenza $\langle a_1, \ldots, a_n \rangle$, una \textbf{sottosequenza} è un
sottoinsieme ordinato $\langle a_{i_1}, a_{i_2}, \ldots, a_{i_k} \rangle$ dove la
sequenza di indici $\langle i_1, i_2, \ldots, i_k \rangle$ è strettamente crescente.
\end{definizione}

\begin{esempio}[Sottosequenza]
Data $X = \langle A, B, C, B, D, A, B \rangle$, la sequenza
$Z = \langle B, C, D, B \rangle$ è sottosequenza di $X$ perché i suoi indici in $X$
formano la sequenza crescente $\langle 2, 3, 5, 7 \rangle$.
\end{esempio}

\begin{definizione}[Sottosequenza comune — LCS]
Date due sequenze $X$ e $Y$, una sequenza $Z$ è \textbf{sottosequenza comune} di $X$
e $Y$ se è sottosequenza di entrambe. La \textbf{Massima Sottosequenza Comune}
(\emph{Longest Common Subsequence}, LCS) è la sottosequenza comune di lunghezza
massima.
\end{definizione}

\begin{esempio}[LCS (CLRS)]
Sia $X = \langle A, B, C, B, D, A, B \rangle$ e $Y = \langle B, D, C, A, B, A \rangle$.
La sequenza $\langle B, C, A \rangle$ è una sottosequenza comune; la sottosequenza
comune più lunga è $\langle B, C, B, A \rangle$ (lunghezza 4).
\end{esempio}

\paragraph{Problema formale.}
\textbf{Input:} due sequenze $X$ e $Y$ sullo stesso alfabeto.
\textbf{Output:} la sottosequenza comune di lunghezza massima.

\subsubsection{Fase 1 — Sottostruttura ottima}

Sia $X_i = \langle x_1, \ldots, x_i \rangle$ l'$i$-esimo prefisso di $X$.

\begin{teorema}[Sottostruttura ottima della LCS]
Dati $X = \langle x_1, \ldots, x_m \rangle$, $Y = \langle y_1, \ldots, y_n \rangle$
e una LCS $Z = \langle z_1, \ldots, z_k \rangle$:
\begin{enumerate}
	\item Se $x_m = y_n$, allora $z_k = x_m$ e $Z_{k-1}$ è una LCS di $X_{m-1}$ e
	$Y_{n-1}$.
	\item Se $x_m \neq y_n$ e $z_k \neq x_m$, allora $Z$ è una LCS di $X_{m-1}$ e $Y$.
	\item Se $x_m \neq y_n$ e $z_k \neq y_n$, allora $Z$ è una LCS di $X$ e $Y_{n-1}$.
\end{enumerate}
\end{teorema}

\begin{proof}
\textbf{Punto 1.} Se $x_m = y_n$ ma $Z$ non termina con tale carattere, si può
sostituire l'ultimo elemento di $Z$ con $x_m$ ottenendo una LCS che termina con
$x_m = y_n$. Tolto l'ultimo carattere, la parte rimanente è una LCS di $X_{m-1}$ e
$Y_{n-1}$.

\textbf{Punto 2.} Se $z_k \neq x_m$, allora $Z$ è una LCS di $X_{m-1}$ e $Y$: se
esistesse una sottosequenza comune $W$ di $X_{m-1}$ e $Y$ con $|W| > k$, allora $W$
sarebbe anche sottosequenza comune di $X$ e $Y$, contraddicendo l'ipotesi.

\textbf{Punto 3.} Dimostrazione simmetrica al punto 2.
\end{proof}

\subsubsection{Fase 2 — Definizione ricorsiva del valore}

Sia $c[i,j]$ la lunghezza di una LCS di $X_i$ e $Y_j$.
La soluzione del problema originale è $c[m, n]$.
La ricorrenza è:

\[
c[i,j] =
\begin{cases}
0 & \text{se } i = 0 \text{ o } j = 0 \\
c[i-1,\, j-1] + 1 & \text{se } i,j > 0 \text{ e } x_i = y_j \\
\max\{c[i,\, j-1],\; c[i-1,\, j]\} & \text{se } i,j > 0 \text{ e } x_i \neq y_j
\end{cases}
\]

\begin{hardnote}
Perché non basta il Divide et Impera? Con $i = j = n$, il numero di chiamate
ricorsive nel caso peggiore vale $T(i,j) = 2\binom{i+j}{i} - 1$, valore
\textbf{esponenziale} quando $i$ e $j$ crescono insieme. I sottoproblemi distinti
sono invece soltanto $\Theta(nm)$, uno per ogni coppia di prefissi $(X_i, Y_j)$.
\end{hardnote}

\subsubsection{Fase 3 — Calcolo bottom-up: algoritmo \textsc{LCS-Length}}

\begin{algorithm}[H]
\caption{\textsc{LCS-Length}$(X, Y)$}
\KwIn{$X = \langle x_1,\ldots,x_m\rangle$, $Y = \langle y_1,\ldots,y_n\rangle$}
\KwOut{Tabelle $c[0{:}m,\,0{:}n]$ e $b[1{:}m,\,1{:}n]$}

Siano $b[1{:}m,\,1{:}n]$ e $c[0{:}m,\,0{:}n]$ due tabelle inizializzate a $0$
\tcp*{$b$ serve per ricostruire la LCS}

\For{$i \leftarrow 1$ \KwTo $m$}{
	\For{$j \leftarrow 1$ \KwTo $n$}{
		\eIf{$x_i = y_j$}{
			$c[i,j] \leftarrow c[i-1,j-1] + 1$\;
			$b[i,j] \leftarrow \text{``}\nwarrow\text{''}$\;
		}{
			\eIf{$c[i-1,j] \geq c[i,j-1]$}{
				$c[i,j] \leftarrow c[i-1,j]$\;
				$b[i,j] \leftarrow \text{``}\uparrow\text{''}$\;
			}{
				$c[i,j] \leftarrow c[i,j-1]$\;
				$b[i,j] \leftarrow \text{``}\leftarrow\text{''}$\;
			}
		}
	}
}
\Return{$c$, $b$}
\end{algorithm}

Il tempo di computazione è $\Theta(nm)$.

\subsubsection{Fase 4 — Ricostruzione: algoritmo \textsc{Print-LCS}}

\begin{algorithm}[H]
\caption{\textsc{Print-LCS}$(b, X, i, j)$}
\KwIn{Matrice $b$; sequenza $X$; indici correnti $i$, $j$}

\If{$i = 0$ \KwAnd $j = 0$}{\Return}

\eIf{$b[i,j] = \text{``}\nwarrow\text{''}$}{
	\textsc{Print-LCS}$(b, X, i-1, j-1)$\;
	stampa $x_i$\;
}{
	\eIf{$b[i,j] = \text{``}\uparrow\text{''}$}{
		\textsc{Print-LCS}$(b, X, i-1, j)$\;
	}{
		\textsc{Print-LCS}$(b, X, i, j-1)$\;
	}
}
\end{algorithm}

Si invoca come \textsc{Print-LCS}$(b, X, |X|, |Y|)$. Il tempo di computazione è
$O(n + m)$.

\begin{esempio}[LCS su ABCBDAB / BDCABA]
Per $X = \langle A,B,C,B,D,A,B \rangle$ e $Y = \langle B,D,C,A,B,A \rangle$,
l'algoritmo restituisce la LCS $\langle B, C, B, A \rangle$ di lunghezza 4,
leggendo le frecce $\nwarrow$ nella tabella $b$.
\end{esempio}

\begin{figure}[H]
	\centering
	\renewcommand{\arraystretch}{1.1}
	\setlength{\tabcolsep}{3pt}
	\begin{tabular}{c|cccccccc}
		\toprule
		 & $\epsilon$ & $B$ & $D$ & $C$ & $A$ & $B$ & $A$ \\
		\midrule
		$i$ & $\epsilon$ & 0 & 0 & 0 & 0 & 0 & 0 & 0 \\
		$A$ & 0 & 0 & 1 & 1 & 1 & 1 & 1 & 1 \\
		$B$ & 0 & 0 & 1 & 1 & 2 & 2 & 2 & 2 \\
		$C$ & 0 & 0 & 1 & 1 & 2 & 2 & 2 & 3 \\
		$B$ & 0 & 0 & 1 & 1 & 2 & 2 & 3 & 3 \\
		$D$ & 0 & 0 & 1 & 2 & 2 & 2 & 3 & 3 \\
		$A$ & 0 & 0 & 1 & 2 & 2 & 2 & 3 & \cellcolor{cyan!15}4 \\
		$B$ & 0 & 0 & 1 & 2 & 2 & 2 & 3 & 4 \\
		\bottomrule
	\end{tabular}
	\caption{Tabella $c[i,j]$ per LCS di $X=\langle A,B,C,B,D,A,B\rangle$ e $Y=\langle B,D,C,A,B,A\rangle$ (valore ottimo $4$ in $c[7,7]$).}
\end{figure}

% -------------------------------------------------------
\subsection{Massima Sottosequenza Crescente}
% -------------------------------------------------------

\begin{definizione}[Sottosequenza crescente]
Una sottosequenza $a_{i_1}, a_{i_2}, \ldots, a_{i_k}$ di interi è
\textbf{crescente} se $a_{i_1} < a_{i_2} < \cdots < a_{i_k}$.
\end{definizione}

\paragraph{Problema.}
\textbf{Input:} sequenza di interi $a_1, \ldots, a_n$.
\textbf{Output:} sottosequenza crescente di lunghezza massima.

\begin{esempio}[MSC]
La sequenza $5, 2, 8, 6, 3, 6, 9, 7$ ha massima sottosequenza crescente
$2, 3, 6, 9$ (lunghezza 4).
\end{esempio}

\subsubsection{Riduzione a cammino massimo in un DAG}

Si costruisce un DAG $G = (V, E)$ dove esiste l'arco $(i, j)$ se e solo se
$i < j$ e $a_i < a_j$. Trovare la massima sottosequenza crescente è
\textbf{equivalente} a trovare il cammino più lungo nel DAG.

\begin{figure}[H]
	\centering
	\begin{tikzpicture}[
		v/.style={circle, draw=primary, minimum size=6mm, font=\scriptsize},
		e/.style={-{Stealth[length=1.8mm]}, primary, thick}
		]
		\foreach \i/\lab in {1/5,2/2,3/8,4/6,5/3,6/6,7/9,8/7}{
			\node[v] (n\i) at (90-45*\i:2.2cm) {\lab};
		}
		\foreach \i/\j in {1/3,1/7,2/5,2/7,3/7,4/7,5/6,5/7,6/7}{
			\draw[e] (n\i) -- (n\j);
		}
		\draw[e, red!70!black, thick] (n2) -- (n5) -- (n6) -- (n7);
	\end{tikzpicture}
	\caption{DAG per $a=\langle 5,2,8,6,3,6,9,7\rangle$: cammino massimo $2\to 3\to 6\to 9$ (lunghezza 4).}
\end{figure}

\subsubsection{Definizione ricorsiva}

Sia $L(j)$ la lunghezza del cammino più lungo che termina nel nodo $j$:
\[
L(j) = 1 + \max\bigl(\{0\} \cup \{L(i) \mid (i,j) \in E\}\bigr).
\]

Il valore della soluzione ottima è $\max_j L(j)$.

\subsubsection{Algoritmo bottom-up (versione valutativa)}

\begin{algorithm}[H]
\caption{\textsc{MSCresc-V}$(a_1, \ldots, a_n)$}
\KwIn{Sequenza di interi $a_1, \ldots, a_n$}
\KwOut{Lunghezza della massima sottosequenza crescente}

Costruisci il DAG trasposto $G^T = (V, E^T)$ da $(a_1, \ldots, a_n)$\;
$L \leftarrow$ array di $n$ interi \tcp*{soluzioni parziali}
$\mathit{ottimo} \leftarrow -\infty$\;

\For{$j \leftarrow 1$ \KwTo $n$}{
	$\mathit{max} \leftarrow 0$\;
	\For{ogni $i$ tale che $(j, i) \in E^T$}{
		\If{$L[i] > \mathit{max}$}{$\mathit{max} \leftarrow L[i]$}
	}
	$L[j] \leftarrow 1 + \mathit{max}$\;
	\If{$L[j] > \mathit{ottimo}$}{$\mathit{ottimo} \leftarrow L[j]$}
}
\Return{$\mathit{ottimo}$}
\end{algorithm}

\paragraph{Complessità.}
La costruzione del grafo trasposto richiede $O(n^2)$. Il calcolo di tutti gli
$L(j)$ richiede $O(|E|)$, con $|E| \leq \frac{n(n-1)}{2}$.
La complessità totale è $O(n^2)$.

\begin{hardnote}
L'approccio Divide et Impera puro genera $2^{j-1}$ chiamate ricorsive per calcolare
$L(j)$ nel caso peggiore (sequenza completamente crescente), rendendo l'algoritmo
esponenziale. La programmazione dinamica risolve ciascuno degli $n$ stati
esattamente una volta.
\end{hardnote}

% -------------------------------------------------------
\subsection{String Matching Approssimato}
% -------------------------------------------------------

\begin{softnote}
Nella ricerca approssimata di un pattern $P$ in un testo $T$ si ammettono tre tipi
di errore: \emph{sostituzione} (caratteri corrispondenti diversi),
\emph{gap in $P$} (un carattere di $T$ non allineato con nessun carattere di $P$),
\emph{gap in $T$} (un carattere di $P$ non allineato con nessun carattere di $T$).
Questo problema generalizza lo string matching esatto ed è strettamente correlato
alla \textbf{distanza di edit} (Levenshtein).
\end{softnote}

\begin{definizione}[Occorrenza $k$-approssimata]
Data una stringa $T$ (testo, lunghezza $n$) e una stringa $P$ (pattern,
lunghezza $m < n$), un'\textbf{occorrenza $k$-approssimata} di $P$ in $T$ è un
allineamento di $P$ con una sottostringa di $T$ in cui compaiono esattamente $k$
errori dei tipi: sostituzione, gap in $P$, gap in $T$.
\end{definizione}

\paragraph{Problema formale.}
\textbf{Input:} testo $T$ di lunghezza $n$, pattern $P$ di lunghezza $m < n$.
\textbf{Output:} occorrenza $k$-approssimata minima (rispetto a $k$) di $P$ in $T$.

\subsubsection{Definizione ricorsiva}

Sia $K[i,j]$ il minimo numero di errori nell'approssimare $P[1,\ldots,i]$ con una
sottostringa di $T[1,\ldots,j]$ che termina in posizione $j$. La soluzione ottima è
$k = \min_j K[m, j]$.

Equazione di ricorrenza unificata:
\[
K[i,j] = \min\bigl\{\,\mathrm{diff}(i,j) + K[i-1,j-1],\;\;
1 + K[i-1,j],\;\;
1 + K[i,j-1]\,\bigr\}
\]
dove $\mathrm{diff}(i,j) = 0$ se $P[i] = T[j]$, altrimenti $1$.

\paragraph{Casi base.}
\begin{itemize}
\item $K[i, 0] = i$ \quad (il pattern ha $i$ caratteri in più della stringa vuota);
\item $K[0, j] = 0$ \quad (il pattern vuoto si allinea a costo $0$ a qualunque prefisso di $T$).
\end{itemize}

\subsubsection{Algoritmo \textsc{SMA}}

\begin{algorithm}[H]
\caption{\textsc{SMA}$(T, P)$}
\KwIn{Testo $T$, pattern $P$}
\KwOut{(numero minimo di errori; posizione ultimo carattere dell'occorrenza in $T$)}

$n \leftarrow |T|$;\; $m \leftarrow |P|$\;
\For{$i \leftarrow 0$ \KwTo $n$}{$K[0,i] \leftarrow 0$}
\For{$i \leftarrow 0$ \KwTo $m$}{$K[i,0] \leftarrow i$}

\For{$i \leftarrow 1$ \KwTo $m$}{
	\For{$j \leftarrow 1$ \KwTo $n$}{
		$\mathit{diff} \leftarrow (T[j] = P[i])\ ?\ 0 : 1$\;
		$K[i,j] \leftarrow \min\{\,\mathit{diff} + K[i-1,j-1],\;\; 1+K[i-1,j],\;\; 1+K[i,j-1]\,\}$\;
	}
}

$\mathit{Min} \leftarrow +\infty$;\; $\mathit{pos} \leftarrow -1$\;
\For{$j \leftarrow 1$ \KwTo $n$}{
	\If{$K[m,j] < \mathit{Min}$}{
		$\mathit{Min} \leftarrow K[m,j]$;\; $\mathit{pos} \leftarrow j$\;
	}
}
\Return{$(\mathit{Min},\; \mathit{pos})$}
\end{algorithm}

La complessità è $\Theta(nm)$.

\begin{esempio}[SMA su ``It's snowy'' / ``sunny'']
Con $T = \texttt{It's snowy}$ e $P = \texttt{sunny}$, la tabella $K$ produce
$\min_j K[5,j] = 3$, confermando che il pattern si abbina al testo con 3 errori.
\end{esempio}

\begin{figure}[H]
	\centering
	\footnotesize
	\renewcommand{\arraystretch}{1.05}
	\begin{tabular}{c|ccccccccccc}
		\toprule
		& $\epsilon$ & I & t & ' & s & \texttt{ } & s & n & o & w & y \\
		\midrule
		$\epsilon$ & 0 & 0 & 0 & 0 & 0 & 0 & 0 & 0 & 0 & 0 & 0 \\
		s & 1 & 1 & 1 & 1 & 0 & 1 & 0 & 1 & 1 & 1 & 1 \\
		u & 2 & 2 & 2 & 2 & 1 & 1 & 1 & 1 & 2 & 2 & 2 \\
		n & 3 & 3 & 3 & 3 & 2 & 2 & 2 & 1 & 2 & 3 & 3 \\
		n & 4 & 4 & 4 & 4 & 3 & 3 & 3 & 2 & 2 & 3 & 4 \\
		y & 5 & 5 & 5 & 5 & 4 & 4 & 4 & 3 & 3 & 3 & \cellcolor{cyan!15}3 \\
		\bottomrule
	\end{tabular}
	\caption{Tabella $K[i,j]$ per \textsc{SMA} con $T=\texttt{It's snowy}$ e $P=\texttt{sunny}$: $\min_j K[5,j]=3$ (evidenziato).}
\end{figure}

% -------------------------------------------------------
\subsection{Problema dello Zaino (Knapsack 0-1)}
% -------------------------------------------------------

\paragraph{Problema formale.}
\textbf{Input:} $n$ oggetti con valore $v_i$ e peso $p_i$ (interi positivi);
portata massima $P$.
\textbf{Output:} vettore $x \in \{0,1\}^n$ che massimizza
$\sum_{i=1}^n x_i v_i$ con vincolo $\sum_{i=1}^n x_i p_i \leq P$.

\begin{nota}
Questa versione è chiamata \textbf{Zaino 0-1} perché ogni oggetto o si prende
($ x_i = 1$) o si lascia ($x_i = 0$). Esistono anche le versioni \emph{frazionaria}
(risolvibile greedy) e \emph{con ripetizione}.
\end{nota}

\subsubsection{Sottostruttura ottima}

Si considera la soluzione ottima per i primi $j$ oggetti con capacità $w$.
\begin{itemize}
\item Se l'oggetto $j$ \textbf{non è} nella soluzione: la soluzione è ottima per
i primi $j-1$ oggetti con capacità $w$.
\item Se l'oggetto $j$ \textbf{è} nella soluzione: toltolo, rimane una soluzione
ottima per i primi $j-1$ oggetti con capacità $w - p_j$.
\end{itemize}
In entrambi i casi: \emph{taglia e incolla} mostra che la soluzione ottima del
problema originale contiene la soluzione ottima del sottoproblema.

\subsubsection{Definizione ricorsiva}

Sia $K[j,w]$ il massimo valore ottenibile usando uno zaino di capacità $w$ e i
primi $j$ oggetti. La soluzione è $K[n, P]$.

\[
K[j,w] =
\begin{cases}
0 & \text{se } j = 0 \text{ o } w = 0 \\
K[j-1,\, w] & \text{se } p_j > w \\
\max\!\bigl\{K[j-1,\, w],\; K[j-1,\, w-p_j]+v_j\bigr\} & \text{altrimenti}
\end{cases}
\]

\subsubsection{Algoritmo \textsc{Zaino}}

\begin{algorithm}[H]
\caption{\textsc{Zaino}$(P, p, v)$}
\KwIn{Portata $P$; array pesi $p$; array valori $v$}
\KwOut{Valore ottimo dello zaino}

$n \leftarrow |v|$\;
\tcp*{In Java le tabelle sono già inizializzate a 0}
\For{$j \leftarrow 1$ \KwTo $n$}{
	\For{$w \leftarrow 1$ \KwTo $P$}{
		\eIf{$p[j] > w$}{
			$K[j,w] \leftarrow K[j-1,w]$ \tcp*{oggetto $j$ troppo pesante}
		}{
			$K[j,w] \leftarrow \max\bigl(K[j-1,w],\; K[j-1,w-p[j]]+v[j]\bigr)$\;
		}
	}
}
\Return{$K[n,P]$}
\end{algorithm}

\paragraph{Complessità.} La complessità è $\Theta(nP)$.

\begin{hardnote}
$\Theta(nP)$ è \textbf{pseudo-polinomiale}, non polinomiale nella dimensione
dell'input. Se $k$ è il numero di bit per rappresentare $P$, allora
$k = \Theta(\log P)$, cioè $P = \Theta(2^k)$: il costo può essere esponenziale
nella lunghezza della rappresentazione binaria di $P$.
Il problema Zaino 0-1 è NP-difficile.
\end{hardnote}

\begin{esempio}[Zaino 0-1]
Con $v = [6, 13, 4, 4]$, $p = [3, 13, 5, 3]$, $P = 13$, lo zaino ottimo contiene
gli oggetti $\{x_1, x_3, x_4\}$ con valore totale $6+4+4=14$ e peso $3+5+3=11$.
\end{esempio}

\begin{figure}[H]
	\centering
	\footnotesize
	\renewcommand{\arraystretch}{1.05}
	\begin{tabular}{c|*{14}{c}}
		\toprule
		$w$ & 0 & 1 & 2 & 3 & 4 & 5 & 6 & 7 & 8 & 9 & 10 & 11 & 12 & 13 \\
		\midrule
		$j=0$ & 0 & 0 & 0 & 0 & 0 & 0 & 0 & 0 & 0 & 0 & 0 & 0 & 0 & 0 \\
		$j=1$ & 0 & 0 & 0 & 6 & 6 & 6 & 6 & 6 & 6 & 6 & 6 & 6 & 6 & 6 \\
		$j=2$ & 0 & 0 & 0 & 6 & 6 & 6 & 6 & 6 & 6 & 6 & 6 & 6 & 6 & 6 \\
		$j=3$ & 0 & 0 & 0 & 6 & 6 & 6 & 6 & 6 & 6 & 6 & 6 & 6 & 10 & 10 \\
		$j=4$ & 0 & 0 & 0 & 6 & 6 & 6 & 6 & 6 & 6 & 6 & 6 & 6 & 10 & \cellcolor{cyan!15}14 \\
		\bottomrule
	\end{tabular}
	\caption{Tabella $K[j,w]$ per lo zaino $v=[6,13,4,4]$, $p=[3,13,5,3]$, $P=13$ (valore ottimo $14$).}
\end{figure}

% -------------------------------------------------------
\subsection{Casi comuni di sottoproblemi}
% -------------------------------------------------------

\begin{nota}
Caratterizzare correttamente la decomposizione in sottoproblemi richiede
creatività ed esperienza. Quattro schemi ricorrono più frequentemente:
\begin{enumerate}
	\item \textbf{Prefisso singolo.} Input $x_1,\ldots,x_n$; il sottoproblema
	considera $x_1,\ldots,x_i$. Numero di stati: $O(n)$.
	\item \textbf{Doppio prefisso.} Input $x_1,\ldots,x_m$ e $y_1,\ldots,y_n$;
	il sottoproblema considera $x_1,\ldots,x_i$ e $y_1,\ldots,y_j$.
	Numero di stati: $O(nm)$.
	\item \textbf{Intervallo.} Input $x_1,\ldots,x_n$; il sottoproblema
	considera $x_i,\ldots,x_j$. Numero di stati: $O(n^2)$
	(tabella $n \times n$ triangolare superiore).
	\item \textbf{Sottoalbero.} Input: albero radicato di $n$ nodi; il
	sottoproblema è un sottoalbero radicato. Numero di stati: $O(n)$.
\end{enumerate}
\end{nota}

% -------------------------------------------------------
\subsection{Memoization}
% -------------------------------------------------------

\begin{definizione}[Memoization]
Tecnica (termine coniato da Donald Michie nel 1968, senza la ``r'' di
\emph{memorization}) che, in un programma ricorsivo, memorizza i risultati delle
chiamate già eseguite in una struttura di supporto (hash table, array, matrice)
indicizzata dai parametri, per riutilizzarli anziché ricalcolarli.
\end{definizione}

\begin{softnote}
\textbf{Confronto tra i tre approcci:}
\begin{itemize}
	\item \textbf{Tabulazione bottom-up}: risolve \emph{tutte} le sottoistanze
	rispettando l'ordine di dipendenza. Iterativa, costanti moltiplicative
	ridotte.
	\item \textbf{Divide et impera}: risolve solo le sottoistanze necessarie ma
	ricalcola quelle comuni ogni volta.
	\item \textbf{Memoization}: mantiene la struttura ricorsiva del divide et
	impera ma riusa i risultati già incontrati. Può essere più efficiente
	della tabulazione se molte sottoistanze non sono necessarie; in generale
	ha overhead maggiore per le chiamate ricorsive.
\end{itemize}
\end{softnote}

\subsubsection{Esempio: Zaino con memoization}

\begin{algorithm}[H]
\caption{\textsc{Zaino-Memo}$(i, w)$}
\KwIn{Indice oggetto $i$, capacità residua $w$; $P$, $v$, $p$ globali}
\KwOut{Valore ottimo per i primi $i$ oggetti con capacità $w$}

\If{$K$.\textsc{containsKey}$(i, w)$}{\Return{$K$.\textsc{get}$(i, w)$}}
\If{$i = 0$ \KwAnd $w = 0$}{\Return{$0$}}
\eIf{$p[i] > w$}{
	$k \leftarrow$ \textsc{Zaino-Memo}$(i-1, w)$ \tcp*{troppo pesante}
}{
	$k \leftarrow \max\bigl($\textsc{Zaino-Memo}$(i-1,w)$,\;
	\textsc{Zaino-Memo}$(i-1,w-p[i])+v[i]\bigr)$\;
}
$K$.\textsc{put}$((i,w),\, k)$\;
\Return{$k$}
\end{algorithm}

Complessità: $O(nP)$, identica alla versione tabulare, ma con costanti
moltiplicative più grandi per l'overhead delle chiamate ricorsive.